\documentclass{article}

\IfFileExists{styles/main.sty}{%
  \usepackage[dstu]{styles/main}%
}{%
  \usepackage[dstu]{../../styles/main}%
}

\usepackage{xurl}
\usepackage{cleveref}

\IfFileExists{references.bib}{%
  \addbibresource{references.bib}%
}{%
  \addbibresource{pws/3/references.bib}%
}

\IfFileExists{figures/}{%
  \graphicspath{{./figures/}}%
}{%
  \graphicspath{{./pws/3/figures/}}%
}

\DefineBibliographyStrings{english}{%
  urlseen = {accessed:},
}

\begin{document}

% ==================== TITLE PAGE ====================
\IfFileExists{docs/title.pdf}{%
  \includepdf[pages=-]{docs/title.pdf}%
}{%
  \includepdf[pages=-]{pws/3/docs/title.pdf}%
}

\selectlanguage{ukrainian}

\crefname{figure}{рис.}{рис.}
\Crefname{figure}{Рисунок}{Рисунки}
\crefname{table}{табл.}{табл.}
\Crefname{table}{Таблиця}{Таблиці}
\crefname{section}{розд.}{розд.}
\Crefname{section}{Розділ}{Розділи}
\crefname{equation}{формула}{формули}
\Crefname{equation}{Формула}{Формули}
\renewcommand{\crefrangeconjunction}{~--~}
\renewcommand{\crefpairconjunction}{ та~}
\renewcommand{\crefmiddleconjunction}{, }
\renewcommand{\creflastconjunction}{ та~}

\setlength{\fboxsep}{0pt}
\setlength{\fboxrule}{0.5pt}

\titleformat{\section}[block]
{\filcenter\bfseries\fontsize{16pt}{24pt}\selectfont}
{}
{0pt}
{\MakeUppercase}
\titlespacing*{\section}{0pt}{20pt}{6pt}

% ==================== CONTENT ====================
\section{ПРАКТИЧНА РОБОТА № 3}

\textbf{Тема:} Пошук наукової літератури за допомогою спеціалізованих пошукових систем та бібліографічних баз.

\textbf{Мета:} Навчитися здійснювати пошук наукової літератури за допомогою спеціалізованих пошукових систем та бібліографічних баз (Scopus, Google Академія) релевантних темі дипломної роботи.

\subsection{Постановка завдання}
Здійснити пошук за ключовими словами релевантних до теми дипломної роботи наукових праць у системі Google Академія та наукометричній базі Scopus. Під час пошуку у базі Scopus, сформувати необхідний запит, який забезпечить можливість отримання невеликої кількості (до 50) релевантних до теми диплому наукових праць~\cite{izonin2026}.

\subsection{Порядок виконання роботи}

\subsubsection{Тема, мета, об'єкт і предмет дослідження БКР}

\textbf{Тема:} «Розробка системи керування рухом нижніх кінцівок двоногого робота».

\textbf{Мета:} розробити та оцінити систему керування на основі навчання з підкріпленням та інтеграції зовнішніх сенсорів для забезпечення ходи двоногого робота по нерівних і неструктурованих поверхнях.

\textbf{Об'єкт:} локомоція двоногого робота на нерівних і неструктурованих поверхнях.

\textbf{Предмет:} системи керування на основі навчання з підкріпленням, що інтегрують екстероцепцію для стійкого пересування на нерівних і неструктурованих поверхнях.

Ключові слова БКР описано у~\cref{tab:keywords}.

\begin{table}[H]
  \caption{Ключові слова українською та англійською мовами.}
  \label{tab:keywords}
  \centering
  \begin{tabular}{|c|l|l|}
    \hline
    \textbf{№} & \textbf{Українська}              & \textbf{Англійська}           \\
    \hline
    1          & гуманоїд                         & humanoid                      \\
    \hline
    2          & локомоція                        & locomotion                    \\
    \hline
    3          & навчання з підкріпленням         & reinforcement learning        \\
    \hline
    4          & мапа висот                       & height map                    \\
    \hline
    5          & симуляція                        & simulation                    \\
    \hline
  \end{tabular}
\end{table}

\subsubsection{Пошук у Google Scholar}

За ключовими словами з~\cref{tab:keywords} спочатку сформовано пошуковий запит українською мовою, однак жодного релевантного результату не знайдено. Потім той самий набір ключових слів використано англійською мовою; з обмеженням «за останній рік» та виключенням патентів і цитувань отримано 107 джерел (\cref{fig:google_scholar_search}).

\begin{figure}[h]
  \centering
  \fbox{\includegraphics[width=0.88\linewidth]{google_scholar.jpg}}
  \caption{Результати пошуку в Google Scholar за англомовним запитом.}
  \label{fig:google_scholar_search}
\end{figure}

До бібліотеки Zotero було додано три найбільш цитовані зі знайдених статті:
\begin{itemize}
  \item Song H. et al. \href{https://doi.org/10.1109/LRA.2026.3664167}{Gait-Adaptive Perceptive Humanoid Locomotion With Real-Time Under-Base Terrain Reconstruction}. 2026~\cite{songGaitAdaptivePerceptiveHumanoid2026};
  \item Sun J. et al. \href{https://doi.org/10.1109/LRA.2026.3703584}{DPL: Depth-Only Perceptive Humanoid Locomotion via Realistic Depth Synthesis and Cross-Attention Terrain Reconstruction}. 2026~\cite{sunDPLDepthOnlyPerceptive2026};
  \item Uddin M. S. \href{https://doi.org/10.1038/s41598-025-34956-7}{Adaptive Motion Planning for Legged Robots in Unstructured Terrain Using Deep Reinforcement Learning}. 2026~\cite{uddinAdaptiveMotionPlanning2026}.
\end{itemize}

\subsubsection{Пошук у Scopus}

Спочатку в базі Scopus виконано пошук з використанням усіх англомовних ключових слів з~\cref{tab:keywords}, однак отримано 0 результатів. Після цього ключове слово «height map» виключено з запиту та застосовано такі фільтри:
\begin{itemize}
  \item роки публікації~--- із 2024 по 2026;
  \item лише відкритий доступ;
  \item тип документа~--- лише статті;
  \item мова~--- англійська;
  \item предметна область~--- комп'ютерні науки.
\end{itemize}

Знайдено 14 джерел (\cref{fig:scopus_search}) й отримано такий пошуковий запит: (TITLE-ABS-KEY(humanoid AND locomotion AND "reinforcement learning" AND simulation) AND ( LIMIT-TO ( SUBJAREA, "COMP" ) ) AND ( LIMIT-TO ( DOCTYPE, "ar" ) ) AND ( LIMIT-TO ( LANGUAGE, "English" ) ) AND ( LIMIT-TO ( OA, "all" ) ) AND PUBYEAR > 2023 AND PUBYEAR < 2028)).

\begin{figure}[h]
  \centering
  \fbox{\includegraphics[width=0.88\linewidth]{scopus.jpg}}
  \caption{Результати пошуку в Scopus за уточненим запитом.}
  \label{fig:scopus_search}
\end{figure}

До бібліотеки Zotero додано три найбільш цитовані зі знайдених статті:
\begin{itemize}
  \item Li Z. et al. \href{https://doi.org/10.1177/02783649241285161}{Reinforcement Learning for Versatile, Dynamic, and Robust Bipedal Locomotion Control}. 2025~\cite{liReinforcementLearningVersatile2025};
  \item Li G. et al. \href{https://doi.org/10.1109/LRA.2024.3388842}{AI-CPG: Adaptive Imitated Central Pattern Generators for Bipedal Locomotion Learned Through Reinforced Reflex Neural Networks}. 2024~\cite{liAICPGAdaptiveImitated2024};
  \item Sun S. et al. \href{https://doi.org/10.1016/j.birob.2024.100187}{Leveraging Large Language Models for Comprehensive Locomotion Control in Humanoid Robots Design}. 2024~\cite{sunLeveragingLargeLanguage2024}.
\end{itemize}

\subsubsection{Бібліотека в Zotero}

Усі знайдені під час пошуку джерела додано до тематичної колекції «Science» у бібліографічному менеджері Zotero (\cref{fig:zotero_science}).

\begin{figure}[h]
  \centering
  \fbox{\includegraphics[width=0.88\linewidth]{zotero.jpg}}
  \caption{Колекція «Science» у бібліотеці Zotero з імпортованими джерелами.}
  \label{fig:zotero_science}
\end{figure}

Отримані джерела експортовано у форматі BibTeX, і за допомогою відповідного пакету \LaTeX{} сформовано бібліографію у форматі IEEE (\cref{fig:ieee_bibliography}).

\begin{figure}[h]
  \centering
  \fbox{\includegraphics[width=0.44\linewidth]{ieee.jpg}}
  \caption{Список використаних джерел, оформлений у стилі IEEE.}
  \label{fig:ieee_bibliography}
\end{figure}

\subsection{Висновки}

Під час виконання практичної роботи було вказано тему, мету, предмет і об'єкт дослідження бакалаврської кваліфікаційної роботи, записано ключові слова роботи українською та англійською мовами й виконано пошук наукових робіт за поданою тематикою.

Пошук здійснено за заданими ключовими словами у Google Scholar і Scopus із використанням необхідних фільтрів для обмеження кількості отриманих результатів. У випадку Scopus пошуковий запит експортовано у текстовому форматі, що забезпечує його відтворюваність.

Отриманий практичний досвід пошуку наукових джерел знадобиться при написанні магістерської кваліфікаційної роботи, а саме для розділу з аналізом літературних джерел.

Звіт було створено із використанням \LaTeX{}. Увесь код доступний у публічному репозиторії~\cite{github_science}.

% ==================== BIBLIOGRAPHY ====================
\clearpage
\printbibliography[heading=bibintoc,title={СПИСОК ВИКОРИСТАНИХ ДЖЕРЕЛ}]

\end{document}
